% Generated draft: requires mathematical and source review.
\documentclass[11pt]{article}
\usepackage[margin=25mm]{geometry}
\usepackage{fontspec}
\setmainfont{Latin Modern Roman}
\usepackage{amsmath,amssymb,mathtools,unicode-math}
\setmathfont{Latin Modern Math}
\usepackage{enumitem,needspace,xurl,hyperref}
\setlength{\parindent}{0pt}
\setlength{\parskip}{5pt}
\setlength{\emergencystretch}{3em}
\newcommand{\N}{\mathbb{N}}
\newcommand{\Z}{\mathbb{Z}}
\newcommand{\Q}{\mathbb{Q}}
\newcommand{\R}{\mathbb{R}}
\newcommand{\C}{\mathbb{C}}
\newcommand{\F}{\mathbb{F}}
\newcommand{\A}{\mathbb{A}}
\newcommand{\PP}{\mathbb{P}}
\newcommand{\E}{\mathbb{E}}
\newcommand{\eps}{\varepsilon}
\newcommand{\dd}{\,\mathrm{d}}
\begin{document}
\section*{UnsolvedMath ID: 700674}
\textbf{Source problem code:} catalogue entry
\subsection{Definitions and mathematical statement}
Let \(p\) be a prime and let \(G\) be a group. For \(\varphi\in\operatorname{Aut}(G)\), write \(g^{\varphi^i}=\varphi^i(g)\). An automorphism \(\varphi\) of order exactly \(p\) is called splitting if
\[
\prod_{i=0}^{p-1}g^{\varphi^i}=1\qquad\text{for every }g\in G,
\]
where the factors are multiplied in increasing order of \(i\). For each \(g\in G\), define
\[
H_g=\langle g,g^\varphi,\ldots,g^{\varphi^{p-1}}\rangle.
\]
For a group \(H\), put \(\gamma_1(H)=H\) and \(\gamma_{n+1}(H)=[\gamma_n(H),H]\), where \([A,B]\) is generated by the commutators \(a^{-1}b^{-1}ab\) with \(a\in A\) and \(b\in B\). The group \(H\) is nilpotent if \(\gamma_{c+1}(H)=\{1\}\) for some integer \(c\geq1\). The group \(G\) is locally nilpotent if every subgroup generated by a finite subset of \(G\) is nilpotent.

The supplied source asks: if \(G\) is a \(p'\)-group admitting a splitting automorphism \(\varphi\) of order \(p\), and
\[
\forall g\in G\ \exists c_g\geq1:\quad \gamma_{c_g+1}(H_g)=\{1\},
\]
must it follow that
\[
\forall\text{ finite }S\subseteq G\ \exists c_S\geq1:\quad
\gamma_{c_S+1}(\langle S\rangle)=\{1\}?
\]
No uniform bound on the nilpotency classes of the groups \(H_g\) is assumed.

The source does not define \(p'\)-group. Under the periodic convention, this means that every element of \(G\) has finite order relatively prime to \(p\). A broader possible convention requires only that \(G\) have no element of order \(p\), allowing elements of infinite order. These give different domains for the displayed question. The intended convention requires source confirmation; the supplied record does not resolve it. In particular, no finiteness assumption on \(G\) is supplied.
\subsection{Short English statement}
For a prime \(p\), suppose a \(p'\)-group has a splitting automorphism of order \(p\), and the subgroup generated by each individual element together with its automorphism orbit is nilpotent. Must every finitely generated subgroup be nilpotent? The source leaves the convention for \(p'\)-group unstated.
\subsection{Sources}
\begin{itemize}
\item Local UnsolvedMath registry, supplied complete source record: permanent ID \(700674\), exact source problem code \texttt{catalogue entry}, title \emph{Kourovka Notebook Problem 10.59}. Upstream revision: \texttt{d371ee917ac9647fbb06cd6fc529deec6f3dda82}.
\item Attributed to A. I. Sozutov, proposed in 1986. The supplied background identifies \emph{The Kourovka Notebook (Unsolved Problems in Group Theory)}, 10th issue (1986), original problem number 10.59, and the 21st issue edition (2026), identified as \texttt{arXiv:1401.0300v48}.
\item The supplied collection description names E. I. Khukhro and V. D. Mazurov as editors of the Kourovka Notebook collection covering Issues 1--21 (1965--2026).
\item The registry labels the problem \emph{open}. This is a supplied source label, not an independently verified assessment of the literature. No separate literature-verification date or literature-source links are supplied.
\end{itemize}
\end{document}
